\section{Phase 1: Durable Workflow-Based Counter -- Design, Execution Semantics, and Experimental Evaluation}

\subsection{Introduction}

Serverless computing platforms have evolved to support increasingly complex workflows beyond stateless function execution. Durable execution frameworks enable the definition of long-running, fault-tolerant workflows by persisting execution state and enabling deterministic replay.

This phase investigates the behavior of AWS Lambda Durable Execution through the implementation of a simple counter application. Despite its simplicity, the counter serves as a representative workload to study execution determinism, retry semantics, failure handling, and concurrency isolation.

\subsection{Objectives}

The main objectives of this phase are:

\begin{itemize}
    \item To implement a durable workflow composed of multiple steps.
    \item To evaluate deterministic execution under normal conditions.
    \item To analyze system behavior under controlled failures.
    \item To observe retry semantics and atomicity guarantees.
    \item To validate isolation properties under concurrent execution.
\end{itemize}

\subsection{System Model}

The system is modeled as a sequence of durable steps executed within a workflow context:

\begin{equation}
W = \{ S_1, S_2, S_3 \}
\end{equation}

where:

\begin{itemize}
    \item $S_1$: Initialize state
    \item $S_2$: Apply operation
    \item $S_3$: Build response
\end{itemize}

Each step is executed through a durable execution runtime that ensures persistence and replay capability.

\subsection{Execution Semantics}

The execution model follows a deterministic workflow pattern:

\begin{itemize}
    \item Each step is executed sequentially.
    \item The output of step $S_i$ is the input of step $S_{i+1}$.
    \item State transitions are only committed upon successful step completion.
    \item Failed steps are retried from the last consistent checkpoint.
\end{itemize}

This implies a state transition function:

\begin{equation}
state_{i+1} = f_i(state_i)
\end{equation}

where $f_i$ represents the transformation applied by step $i$.

\subsection{Implementation}

The durable workflow is implemented using AWS Lambda Durable Execution SDK in Python.

\begin{verbatim}
from aws_durable_execution_sdk_python.context import DurableContext, StepContext, durable_step
from aws_durable_execution_sdk_python.execution import durable_execution


@durable_step
def initialize_counter(step_context: StepContext, initial_state: dict | None) -> dict:
    if initial_state is None:
        initial_state = {
            "value": 0,
            "version": 0
        }
    return initial_state


@durable_step
def apply_counter_operation(step_context: StepContext, payload: dict) -> dict:
    state = payload["state"]
    operation = payload.get("operation", "get_value")
    amount = int(payload.get("amount", 1))
    fail_mode = payload.get("fail_mode", "none")

    if fail_mode == "always":
        raise RuntimeError("Simulated transient failure in apply_counter_operation")

    if operation == "increment":
        state["value"] += amount
        state["version"] += 1
    elif operation == "decrement":
        state["value"] -= amount
        state["version"] += 1
    elif operation == "get_value":
        pass
    else:
        raise ValueError(f"Unsupported operation: {operation}")

    return state


@durable_step
def build_response(step_context: StepContext, state: dict) -> dict:
    return {
        "message": "fase1-counter-durable executed successfully",
        "counter_value": state["value"],
        "version": state["version"]
    }


@durable_execution
def lambda_handler(event, context: DurableContext) -> dict:
    initial_state = event.get("state")
    operation = event.get("operation", "get_value")
    amount = int(event.get("amount", 1))
    fail_mode = event.get("fail_mode", "none")

    state = context.step(initialize_counter(initial_state))

    state = context.step(apply_counter_operation({
        "state": state,
        "operation": operation,
        "amount": amount,
        "fail_mode": fail_mode
    }))

    result = context.step(build_response(state))

    return {
        "statusCode": 200,
        "body": result
    }
\end{verbatim}

\subsection{Experimental Design}

Three categories of experiments were conducted.

\subsubsection{Normal Execution}

Evaluates correctness and determinism of the workflow under standard conditions.

\subsubsection{Controlled Failure}

Introduces deterministic failure using:

\begin{verbatim}
fail_mode = "always"
\end{verbatim}

This forces the system to trigger retry logic and exposes internal recovery mechanisms.

\subsubsection{Concurrent Execution}

Multiple invocations ($N=5$) are executed in rapid succession with identical input to evaluate isolation and concurrency behavior.

\subsection{Results}

\subsubsection{Deterministic Execution}

All normal executions produced identical outputs and followed identical execution traces, confirming deterministic behavior.

\subsubsection{Failure and Retry Behavior}

Controlled failures triggered automatic retries. The system repeatedly executed the failing step until retry limits were reached.

Observed properties:

\begin{itemize}
    \item No partial state updates were committed.
    \item Each retry re-executed the step from the last stable checkpoint.
    \item The workflow preserved consistency across retries.
\end{itemize}

\subsubsection{Atomicity Guarantees}

The system enforces step-level atomicity:

\begin{itemize}
    \item A step either completes fully or has no effect.
    \item State mutations inside failed steps are discarded.
    \item Side effects are not externally visible unless the step succeeds.
\end{itemize}

This behavior is analogous to transactional execution in distributed systems.

\subsubsection{Concurrency Analysis}

Five concurrent executions were performed. Results show:

\begin{itemize}
    \item All executions completed successfully.
    \item Each execution maintained an independent state.
    \item No interference or race conditions were observed.
\end{itemize}

Each execution returned:

\begin{equation}
counter\_value = 1
\end{equation}

\subsubsection{Isolation Properties}

The observed behavior confirms:

\begin{itemize}
    \item Strong isolation between workflow instances.
    \item Absence of shared mutable state.
    \item Independent execution histories.
\end{itemize}

\subsubsection{Performance Observations}

Execution times exhibited two distinct regimes:

\begin{itemize}
    \item Cold start: approximately 3.5 seconds
    \item Warm execution: approximately 0.8--0.9 seconds
\end{itemize}

This indicates container reuse and runtime optimization after initial invocation.

\subsection{Discussion}

The experimental results indicate that AWS Durable Execution provides:

\begin{itemize}
    \item Deterministic orchestration of workflows
    \item Built-in fault tolerance via retry mechanisms
    \item Strong consistency through step-level atomicity
    \item High isolation across concurrent executions
\end{itemize}

However, the system does not implement shared state semantics. Each workflow instance operates independently, which limits its applicability for workloads requiring coordinated state updates.

\subsection{Limitations}

The current implementation has the following limitations:

\begin{itemize}
    \item No shared state across executions
    \item No conflict resolution mechanisms
    \item No support for distributed coordination between workflows
\end{itemize}

\subsection{Threats to Validity}

\begin{itemize}
    \item The workload is simplified (counter abstraction).
    \item Results may vary under higher concurrency levels.
    \item Observations are specific to AWS Lambda Durable Execution.
\end{itemize}

\subsection{Conclusion}

This phase demonstrates that durable workflows provide a robust execution model for fault-tolerant serverless applications. The system ensures deterministic execution, strong atomicity, and isolation across concurrent invocations.

The results highlight the suitability of this model for orchestration tasks, while also exposing limitations in scenarios requiring shared mutable state.

\subsection{Future Work}

Future work will explore:

\begin{itemize}
    \item Shared-state implementations using external storage (e.g., DynamoDB)
    \item Comparison with actor-based models
    \item Evaluation under high-concurrency workloads
    \item Performance benchmarking across different orchestration frameworks
\end{itemize}
